%% ============================================================
%%  SOLUCIÓN — Ejercicios 3.1 y 3.2: Bibliografía con biblatex
%%  Clase 2 — Después de diapositiva "El sistema de bibliografía"
%%  IMPORTANTE: Este archivo requiere compilación especial:
%%  F5 → Tools → Bibliography (F11) → F5 → F5
%%  O en terminal: pdflatex → biber → pdflatex → pdflatex
%% ============================================================

\documentclass[12pt, a4paper]{article}
\usepackage[utf8]{inputenc}
\usepackage[T1]{fontenc}
\usepackage[spanish]{babel}

% --- BIBLIOGRAFÍA CON BIBLATEX ---
\usepackage[style=authoryear, backend=biber, maxnames=3]{biblatex}
\addbibresource{referencias.bib}  % Asegurate de que el archivo .bib exista

\title{Capital humano y retornos a la educación}
\author{Tu Nombre}
\date{\today}

\begin{document}
\maketitle

\section{Introducción}

El modelo de capital humano, introducido por \textcite{mincer1974}, plantea 
que los salarios crecen con la educación y la experiencia acumulada en el 
mercado laboral. Este modelo ha sido extensamente validado empíricamente en 
múltiples contextos y países.

\section{Marco teórico}

Las instituciones coloniales tienen efectos persistentes en el desarrollo 
económico \parencite{acemoglu2001}. Este hallazgo sugiere que los 
determinantes de largo plazo del crecimiento económico no son solo políticas 
de corto plazo, sino decisiones institucionales tomadas hace siglos.

\section{Metodología}

El instrumento de cercanía a universidades fue propuesto originalmente por 
\citeauthor{card1995} (\citeyear{card1995}) para obtener estimadores 
consistentes del retorno a la educación bajo endogeneidad. Este enfoque ha 
sido ampliamente adoptado en la literatura de economía laboral.

\section{Resultados}

De acuerdo con \textcite{world_bank_2023}, los migrantes y refugiados tienen 
efectos económicos significativos en las sociedades que los reciben. El 
reporte del Banco Mundial documenta tanto desafíos como oportunidades.

\section{Implicancias de política}

Combinando los hallazgos de \textcite{mincer1974} con investigaciones más 
recientes, podemos concluir que:

\begin{enumerate}
  \item La inversión en educación tiene altos retornos privados.
  \item Los retornos varían significativamente según contexto y tipo de educación.
  \item Existen barreras importantes para que ciertos grupos accedan a educación 
        de calidad.
\end{enumerate}

\newpage
\printbibliography[title={Referencias}]

\end{document}
